\section{My impact}

\textbf{3a. How my contributions drove the team forward.} The honest thesis of this section is the one I wrote into our group complaint letter: \emph{the simulation framework we built ourselves is the only reason this project produced any verifiable engineering output at all}. I want to defend that claim with mechanism, not assertion, because it is exactly the kind of line that at Merit-level reads as self-flattery.

The anchor is a full Kolb cycle, which I name stage-by-stage because I used it deliberately rather than retrospectively. My \emph{concrete experience} was three collapsed flight days: the 2026-03-11 weather cancellation, the 2026-03-30 hardware session that produced no airborne minutes, and a third field day that ended with our pilot standing on a wet field while Edward and I debugged a corrupt \texttt{calibration\_data.npz} that was mangling every frame. After the first cancellation I noticed, in my \emph{reflective observation} stage, something harder than ``we lost a day'': the team's entire definition of ``progress'' had silently assumed hardware access we did not have, and the rest of the team was now idle because their planned work needed a drone in the air. I was blocked too, but I had a laptop. That asymmetry was the data point. My \emph{abstract conceptualisation} was the one move I am most proud of in this project: I decided that if \texttt{vision.py} exposed a single interface (\texttt{detect\_in\_image} returning \texttt{(found, x, y, conf)}) and auto-selected between Ultralytics on the laptop and TFLite on the Pi, and if \texttt{config.py} auto-detected its hardware, then the same state machine, the same geofence, and the same search pattern could run in both places --- meaning anything that ran in simulation would, in principle, run on hardware. The delta between sim and real would be the thickness of one module, not the thickness of a rewrite. My \emph{active experimentation} was six weeks of building and then hammering on that premise without a physical drone: \texttt{simple\_simulator.py} grew to 2508 lines with GPS drift, camera shake and TFLite-backend simulation; \texttt{main.py} was refactored from 1411 to 754 lines across six modules with 146 pytest tests (commit \texttt{session-2026-03-20}); and the \texttt{tests/flight/0a--4} progressive ladder plus 127 unit tests gave me a way to know, before any field day, which scripts could survive contact with hardware. When we finally ran on the Pi the delta was indeed thin --- a bbox-coordinate bug (commit \texttt{3d62e6a}), a double-scaling fix (\texttt{7eb7cc8}) and a deleted calibration file --- and the system came up. That is the evidence. I now see that the simulation framework was not a backup plan but a \emph{statement}: that verifiable engineering is possible without hardware access, and this has changed how I think about engineering under institutional constraint. Simulation-first is no longer a personal preference for me; it is a claim about what the discipline owes.

Three other contributions matter here and I mention them compactly. I unblocked Robin by building him a dedicated HTTP integration path (\texttt{passive\_watch\_2.py}, commit \texttt{a5b1ab7}) with a parsing-ready JSON schema (\texttt{885cee4}) and the \texttt{--smart-clear} flag he needed to re-lock on a second target (\texttt{cec899e}); he shipped his state machine without touching my CV code. I unblocked Demetro by building his two Final Design Review slides from scratch (\texttt{32fb0c5}) and then writing and polishing his speaking scripts four to five times (\texttt{b2a9cd4}, \texttt{ba667f2}) --- 465 words, 3.1 minutes of tested delivery, which was the difference between his deliverable shipping on time and shipping late. And I drove the group's external communication about the hardware failures: drafting the complaint letter across six tonal iterations (\texttt{26e3255}--\texttt{847518e}) and turning it into a constructive clarification request the whole team signed off on (\texttt{4b0fd91}).

But the sober closing of 3a is that the simulation framework could \emph{not} tell us what it could not tell us. It could not predict motion blur at 5\,m/s with a real IMX296. It could not find the descent stall that we accepted as a SITL artefact and left unresolved. It could not tell us what the GPS timing lag would actually do to a multi-pass lock --- that came out of DJI video analysis, after the fact. The statement my simulation made is true, and the limits of that statement are also true: sim-to-real parity is an argument I can only win about the things I modelled, and I was disciplined about that only because I had been forced to be.

\textbf{3b. What I would add or re-focus on in future.} I mistook technical leadership for team leadership. I now see they are orthogonal --- you can build great tools that do not enable anyone --- and in this project the gap showed up as a metric: the number of times a teammate other than me edited \texttt{simple\_simulator.py} or \texttt{main.py} was zero. The next project must invert this priority: teach first, build second. Three concrete mechanisms, each with a falsifiability test, follow.

First, I will shift from ``building tools I need'' to ``building tools the team needs,'' and each new tool must have at least one named teammate beneficiary documented \emph{before} it ships --- if I cannot name them, I do not start. Second, no file larger than 500 lines will be refactored solo: I will pair on the review, and the observable signal that the mechanism is working is whether the pair partner can edit that file independently within 48 hours of the walkthrough. Third, I will schedule two hours per week of explicit teaching sessions, and I will grade them not by whether I explained something but by whether the teammate completes a handoff task within 48 hours without asking me a second question. These are tests, not wishes; if after four weeks the observable signals are not firing, I have to change the mechanism rather than reassure myself that I meant well.
